%% Voice Coaching System — Beamer presentation
%% Compile: pdflatex presentation.tex  (run twice)
\documentclass[aspectratio=169,9pt]{beamer}

\usepackage{amsmath,amssymb}
\usepackage{booktabs}
\usepackage{array}
\usepackage{xcolor}

%% ── Theme ────────────────────────────────────────────────────────────────────
\usetheme{Madrid}

\definecolor{darkblue}{RGB}{26,58,92}
\definecolor{midblue}{RGB}{31,119,180}
\definecolor{orng}{RGB}{210,100,10}
\definecolor{dkgrn}{RGB}{30,130,30}
\definecolor{lblue}{RGB}{214,230,248}

\setbeamercolor{palette primary}{bg=darkblue,fg=white}
\setbeamercolor{palette secondary}{bg=midblue,fg=white}
\setbeamercolor{palette tertiary}{bg=darkblue,fg=white}
\setbeamercolor{palette quaternary}{bg=darkblue,fg=white}
\setbeamercolor{frametitle}{bg=darkblue,fg=white}
\setbeamercolor{title}{bg=darkblue,fg=white}
\setbeamercolor{structure}{fg=darkblue}
\setbeamercolor{block title}{bg=midblue,fg=white}
\setbeamercolor{block body}{bg=lblue!40,fg=black}
\setbeamercolor{block title alerted}{bg=orng,fg=white}
\setbeamercolor{block body alerted}{bg=orng!12,fg=black}
\setbeamercolor{block title example}{bg=dkgrn,fg=white}
\setbeamercolor{block body example}{bg=dkgrn!10,fg=black}
\setbeamerfont{frametitle}{size=\small,series=\bfseries}
\setbeamerfont{block title}{size=\small,series=\bfseries}
\setbeamertemplate{navigation symbols}{}

%% ── Notation ─────────────────────────────────────────────────────────────────
\newcommand{\fh}{\hat{f}}
\newcommand{\vh}{\hat{v}}
\newcommand{\sh}{\hat{s}}
\newcommand{\eh}{\hat{e}}
\newcommand{\ah}{\hat{a}}
\newcommand{\bh}{\hat{b}}
\newcommand{\dh}{\hat{d}}

%% ── Speaker notes ────────────────────────────────────────────────────────────
%% To show note pages after each slide, uncomment the two lines below and
%% recompile.  Remove them again for the clean presentation version.
%% \usepackage{pgfpages}
%% \setbeameroption{show notes on second screen=right}
\setbeamertemplate{note page}[plain]

%% ── Title ────────────────────────────────────────────────────────────────────
\title[Voice Coaching System]{Voice Coaching System\\for Singing Performance Evaluation}
\subtitle{Reference-Guided Pipeline $\cdot$ Three-Model Architecture $\cdot$ Deterministic Scoring}
\author{}
\institute{Drexel University}
\date{2026}

% ─────────────────────────────────────────────────────────────────────────────
\begin{document}

%% ─── 1. Title ────────────────────────────────────────────────────────────────
\begin{frame}
  \titlepage
  \note{
    Welcome everyone. Today I will be presenting the Voice Coaching System for
    Singing Performance Evaluation, developed at the IMAPLE Lab at Drexel
    University.

    The central idea of this work is to give singers objective, interpretable
    feedback by comparing their recorded audio against a musical reference score.
    We use three separate neural models to extract measurable parameters from the
    audio, and then all of the coaching scores are computed using explicit
    mathematical equations --- the models never directly predict a score.

    I will walk through the tracked parameters, the metrics we calculate, how
    those metrics are turned into scores, the three models we built, and finally
    some results.
  }
\end{frame}

%% ─── 2. System Overview ──────────────────────────────────────────────────────
\begin{frame}{System Overview}
  \note{
    The key design principle is that neural models act as \textit{feature
    extractors}, not score predictors.  Every coaching number comes from an
    explicit equation you can inspect.

    The pipeline has six stages.  We start with a raw WAV recording of the
    singer.  Three models then extract performance-side features: pitch with
    voicing decisions, note onset and offset times, and phoneme boundaries.
    We load ground-truth reference data from two file types: MusicXML gives us
    note pitches, timings, and tempo from the musical score; TextGrid gives us
    word- and phoneme-level timing boundaries from Praat annotations.  We
    compare predictions against the reference to compute objective error
    metrics, normalise those errors to a zero-to-one-hundred scale, and finally
    a rule-based engine produces human-readable strengths and weaknesses.

    The dataset is GTSinger, English alto-1 subset.
  }
  Neural models estimate measurable parameters; all scores are computed by
  \textbf{deterministic equations} — models never directly predict a score.

  \medskip
  \begin{columns}[T]
    \column{0.46\textwidth}
      \begin{block}{Pipeline}
        \small
        \begin{enumerate}\setlength\itemsep{1pt}
          \item \textbf{Audio Input} — 16\,kHz mono WAV
          \item \textbf{Feature Extraction} — pitch, note boundaries, phonemes
          \item \textbf{Reference Data} — MusicXML + TextGrid
          \item \textbf{Metric Computation} — predicted vs.\ reference
          \item \textbf{Scoring} — normalise $\rightarrow$ [0,\,100]
          \item \textbf{Interpretation} — rule-based feedback
        \end{enumerate}
      \end{block}

    \column{0.50\textwidth}
      \begin{block}{Reference Sources}
        \small
        \renewcommand{\arraystretch}{1.15}
        \begin{tabular}{@{}ll@{}}
          \toprule
          Source   & Provides \\
          \midrule
          MusicXML & note pitch, onset, offset, duration, tempo \\
          TextGrid & phoneme \& word boundary times \\
          \bottomrule
        \end{tabular}
        \medskip\par
        \textbf{Dataset:} GTSinger (English, EN-Alto-1)
      \end{block}
  \end{columns}
\end{frame}

%% ─── 3. Parameters We Track ──────────────────────────────────────────────────
\begin{frame}{Parameters We Track}
  \note{
    Before defining any metrics, I want to establish the notation we use
    throughout the rest of the talk.

    On the performance side --- extracted from the singer's audio --- we have
    six quantities.  F-hat-t is the fundamental frequency at each 10-millisecond
    frame.  V-hat-t is a binary voiced or unvoiced flag.  S-hat-n and E-hat-n
    are the predicted onset and offset times for note n.  D-hat-n is simply the
    difference, the predicted duration.  And A-hat-p and B-hat-p are the
    predicted start and end times for phoneme p.

    The voiced set V is the intersection of frames where both the prediction and
    the reference have a positive frequency --- that is where pitch scoring is
    valid.

    On the reference side, everything comes from the score files.  M-n is the
    MIDI pitch from MusicXML, which we convert to a frequency R-n using the
    standard formula.  S-n, E-n, D-n are the reference timings.  A-p and B-p
    are the TextGrid phoneme boundaries.

    Two sign conventions to keep in mind: a positive onset deviation means the
    singer came in late, and a positive cent error means the singer is sharp.
  }
  \begin{columns}[T]
    \column{0.48\textwidth}
      \begin{block}{Performance Side — extracted from WAV}
        \small
        \renewcommand{\arraystretch}{1.25}
        \begin{tabular}{@{}lp{5cm}@{}}
          $\fh_t$ & Frame-level F0 (Hz) at frame $t$ \\
          $\vh_t \in \{0,1\}$ & Voiced/unvoiced flag per frame \\
          $\sh_n$ & Predicted onset time of note $n$ \\
          $\eh_n$ & Predicted offset time of note $n$ \\
          $\dh_n = \eh_n - \sh_n$ & Predicted duration of note $n$ \\
          $\ah_p,\bh_p$ & Predicted phoneme $p$ start \& end \\
        \end{tabular}
        \medskip\par
        \textbf{Voiced set:}\quad $V = \{t : \fh_t > 0,\; r_t > 0\}$
      \end{block}

    \column{0.48\textwidth}
      \begin{block}{Reference Side — from score files}
        \small
        \renewcommand{\arraystretch}{1.25}
        \begin{tabular}{@{}lp{5cm}@{}}
          $m_n$ & MIDI pitch of note $n$ (MusicXML) \\
          $r_n = 440\cdot 2^{(m_n-69)/12}$ & Reference frequency (Hz) \\
          $s_n,\; e_n$ & Reference onset \& offset (s) \\
          $d_n = e_n - s_n$ & Reference duration (s) \\
          $a_p,\; b_p$ & Reference phoneme $p$ start \& end \\
        \end{tabular}
        \medskip\par
        \small
        \textit{Positive onset deviation} $= $ predicted is later.\\
        \textit{Positive cent error} $= $ predicted is sharper.
      \end{block}
  \end{columns}
\end{frame}

%% ─── 4. Metrics — Pitch ──────────────────────────────────────────────────────
\begin{frame}{Metrics --- Pitch}
  \note{
    We compute four pitch metrics.  They all work on the same input: delta-n,
    the pitch deviation in cents for each matched note pair.

    Pitch Accuracy is the simplest --- the fraction of notes whose deviation
    is within plus or minus tau cents, where tau is set to 50 cents, roughly a
    quarter tone.

    MACE, Mean Absolute Cent Error, is the average absolute deviation.  It does
    not require choosing a tolerance, so it gives a continuous sense of how far
    off the singer tends to be.

    Pitch RMSE is the root-mean-square deviation.  Because it squares the errors
    before averaging, it is more sensitive to occasional large jumps than MACE
    is.

    Note-Level Pitch Accuracy is numerically identical to Pitch Accuracy --- it
    is set to the same value in the code.  The difference is that it also
    populates the per-note breakdown field, which lets us give note-by-note
    feedback.

    All deviations come from the NoteAlignmentMatch structure produced after
    the greedy overlap alignment step.
  }
  Let $N$ = matched notes with valid pitch; $\delta_n$ = pitch deviation in cents
  for matched note $n$;\; $\tau = 50$\,¢.

  \begin{columns}[T]
    \column{0.48\textwidth}
      \begin{block}{Pitch Accuracy}
        \[
          \mathrm{PitchAcc}_\tau
            = \frac{1}{N}\sum_{n=1}^{N}\mathbf{1}(|\delta_n|\le\tau)
        \]
        \footnotesize Fraction of notes whose pitch deviation is within $\pm\tau$ cents.
      \end{block}
      \medskip
      \begin{block}{Pitch RMSE}
        \[
          \mathrm{PitchRMSE} = \sqrt{\frac{1}{N}\sum_{n=1}^{N}\delta_n^2}
        \]
        \footnotesize Penalises large deviations more strongly than MACE.
      \end{block}

    \column{0.48\textwidth}
      \begin{block}{Mean Absolute Cent Error (MACE)}
        \[
          \mathrm{MACE} = \frac{1}{N}\sum_{n=1}^{N}|\delta_n|
        \]
        \footnotesize Average absolute pitch error in cents; no threshold required.
      \end{block}
      \medskip
      \begin{alertblock}{Note-Level Pitch Accuracy}
        \[
          \mathrm{NotePitchAcc}_\tau \;=\; \mathrm{PitchAcc}_\tau
        \]
        \footnotesize Same value; per-note detail in \texttt{PitchMetrics.per\_note}.
      \end{alertblock}
  \end{columns}
  \vspace{4pt}
  \footnotesize
  $\delta_n$ is from \texttt{NoteAlignmentMatch.pitch\_deviation\_cents}
  (computed via the median F0 over the note window).
\end{frame}

%% ─── 5. Metrics — Timing ─────────────────────────────────────────────────────
\begin{frame}{Metrics --- Timing}
  \note{
    For timing we define epsilon-n as the onset deviation in milliseconds ---
    positive means the singer was late.  Similarly epsilon-varepsilon-n is
    the offset deviation.

    For onsets we compute three statistics.  The mean captures systematic bias
    --- if the singer consistently comes in early or late.  The MAE captures
    average accuracy.  The standard deviation captures consistency --- whether
    the timing errors are all similar or wildly variable.

    Offset MAE simply asks how early or late the singer releases each note on
    average.

    Timing Accuracy is the fraction of notes started within 50 milliseconds of
    the reference onset.  The 50-millisecond tolerance is set in pipeline.yaml.

    IOI-MAE stands for Inter-Onset Interval Mean Absolute Error.  An
    inter-onset interval is the time between the start of one note and the
    start of the next.  We compare the sequence of predicted intervals against
    the reference intervals.  This captures rhythmic consistency independently
    of any absolute offset --- a singer who is consistently 200 milliseconds
    late but perfectly in rhythm will have low IOI-MAE but high onset MAE.
  }
  Let $\epsilon_n = (\sh_n - s_n)\times 1000$\,ms (onset deviation) and
  $\varepsilon_n = (\eh_n - e_n)\times 1000$\,ms (offset deviation).

  \begin{columns}[T]
    \column{0.48\textwidth}
      \begin{block}{Onset Error}
        \[
          \overline{\epsilon} = \frac{1}{N}\!\sum_n\epsilon_n,\quad
          \mathrm{MAE}_\text{on} = \frac{1}{N}\!\sum_n|\epsilon_n|
        \]
        \[
          \sigma_\text{on} = \sqrt{\tfrac{1}{N}\textstyle\sum_n(\epsilon_n-\overline{\epsilon})^2}
        \]
        \footnotesize Mean (bias), MAE (accuracy), std (consistency) of onset timing.
      \end{block}
      \medskip
      \begin{block}{Offset Error}
        \[
          \mathrm{MAE}_\text{off} = \frac{1}{N}\sum_n|\varepsilon_n|
        \]
        \footnotesize How early or late the singer releases each note.
      \end{block}

    \column{0.48\textwidth}
      \begin{block}{Timing Accuracy}
        \[
          \mathrm{TimingAcc}_\tau
            = \frac{1}{N}\sum_n\mathbf{1}(|\epsilon_n|\le\tau),
            \quad\tau=50\,\text{ms}
        \]
        \footnotesize Fraction of notes entered within 50\,ms of the reference onset.
      \end{block}
      \medskip
      \begin{block}{IOI Mean Absolute Error}
        \[
          \mathrm{IOI\text{-}MAE}
            = \frac{1}{M}\sum_{i=1}^{M}|\Delta\sh^{(i)}-\Delta s^{(i)}|\times 1000
        \]
        \footnotesize $\Delta s^{(i)}\!=\!s_{i+1}\!-\!s_i$;\;
        $M\!=\!\min(N_\text{pred},N_\text{ref})\!-\!1$;\; measures rhythmic consistency.
      \end{block}
  \end{columns}
\end{frame}

%% ─── 6. Metrics — Duration ───────────────────────────────────────────────────
\begin{frame}{Metrics --- Duration}
  \note{
    Duration metrics measure how well the singer holds notes for the correct
    length of time.  A singer can enter a note perfectly on time but still cut
    it too short or hold it too long.

    Duration Error is the mean absolute difference between predicted and
    reference durations in seconds --- a raw error measure.

    Relative Duration Error is more useful because it normalises by the
    reference duration.  An error of 100 milliseconds on a 200-millisecond note
    is a 50 percent error, which is serious.  The same 100 milliseconds on a
    2-second note is only a 5 percent error.

    Duration Ratio is the mean of predicted duration divided by reference
    duration.  A ratio of 1.0 is perfect.  Greater than 1 means the singer
    stretches notes.  Less than 1 means notes are cut short.  For scoring we
    use the absolute deviation of the ratio from 1.

    In the duration category score, relative error has the highest weight at
    60 percent because it is tempo-independent.  Ratio deviation contributes
    20 percent, and phrase consistency --- which measures how consistent the
    errors are across notes --- also contributes 20 percent.
  }
  For matched notes, let $\dh_n = \eh_n - \sh_n$ and $d_n = e_n - s_n$.

  \begin{columns}[T]
    \column{0.48\textwidth}
      \begin{block}{Duration Error}
        \[
          \mathrm{DurErr} = \frac{1}{N}\sum_n|\dh_n - d_n|
        \]
        \footnotesize Absolute difference between predicted and reference duration (s).
      \end{block}
      \medskip
      \begin{block}{Relative Duration Error}
        \[
          \mathrm{RelDurErr} = \frac{1}{N}\sum_n\frac{|\dh_n - d_n|}{d_n}
        \]
        \footnotesize Normalised by reference duration; tempo-independent.
      \end{block}

    \column{0.48\textwidth}
      \begin{block}{Duration Ratio}
        \[
          \mathrm{DurRatio} = \frac{1}{N}\sum_n\frac{\dh_n}{d_n}
        \]
        \footnotesize 1.0\,=\,perfect;\; $>$1\,=\,stretched;\; $<$1\,=\,cut short.
      \end{block}
      \medskip
      \begin{exampleblock}{Scoring (Ratio Component)}
        \[
          \text{scored as }\;|\mathrm{DurRatio} - 1|
        \]
        \footnotesize Deviation from 1.0; combined with RelDurErr into duration score.
      \end{exampleblock}
  \end{columns}
  \vspace{4pt}
  \footnotesize
  Component weights in duration score:\; RelDurErr 0.60,\; Ratio Dev 0.20,\;
  Phrase Consistency 0.20.
\end{frame}

%% ─── 7. Metrics — Lyric ─────────────────────────────────────────────────────
\begin{frame}{Metrics --- Lyric / Phoneme}
  \note{
    Lyric metrics evaluate how well the singer's phoneme and word timing aligns
    with the TextGrid reference annotations.

    Phoneme Boundary Error is the mean absolute onset deviation across all
    matched phoneme pairs, in milliseconds.  Note that we only use the start
    boundary --- not the average of start and end --- because that is what the
    code in lyric\_metrics.py actually computes.

    Word Alignment Accuracy is the fraction of reference words that were
    matched to a predicted word event.  This captures whether the system
    detected the words at all.

    Phoneme Overlap Accuracy checks whether each predicted phoneme segment
    has at least 50 percent temporal overlap with its reference counterpart.
    This is about spatial alignment rather than just onset timing.

    Label Match Rate asks whether the predicted phoneme label actually matches
    the reference label.  This verifies that we predict the right phoneme, not
    just something at roughly the right time.

    The weights in the lyric score are 35 percent for word accuracy, 25 percent
    each for overlap and label match, and 15 percent for boundary timing.
  }
  Let $P$ = matched phoneme pairs;\;
  $\epsilon_p^{(a)} = (\ah_p - a_p)\times 1000$\,ms (phoneme onset deviation).

  \begin{columns}[T]
    \column{0.48\textwidth}
      \begin{block}{Phoneme Boundary Error}
        \[
          \mathrm{PhonBndErr} = \frac{1}{P}\sum_{p=1}^{P}\bigl|\epsilon_p^{(a)}\bigr|
        \]
        \footnotesize Mean absolute onset deviation in ms; onset boundary only.
      \end{block}
      \medskip
      \begin{block}{Word Alignment Accuracy}
        \[
          \mathrm{WordAcc}
            = \frac{|\text{matched words}|}{|\text{reference words}|}
        \]
        \footnotesize Fraction of reference words matched to a predicted word event.
      \end{block}

    \column{0.48\textwidth}
      \begin{block}{Phoneme Overlap Accuracy}
        \[
          \mathrm{OverlapAcc}
            = \frac{1}{P}\sum_p\mathbf{1}(\mathrm{overlap}_p \ge 0.5)
        \]
        \footnotesize Fraction of phoneme pairs with $\ge 50\%$ temporal overlap.
      \end{block}
      \medskip
      \begin{block}{Label Match Rate}
        \[
          \mathrm{LabelMatch}
            = \frac{1}{P}\sum_p\mathbf{1}(\hat{l}_p = l_p)
        \]
        \footnotesize Fraction of matches where predicted phoneme label is correct.
      \end{block}
  \end{columns}
  \vspace{4pt}
  \footnotesize
  Lyric score weights:\; WordAcc 0.35,\; OverlapAcc 0.25,\;
  LabelMatch 0.25,\; PhonBndErr 0.15.
\end{frame}

%% ─── 8. Scoring — Normalization ──────────────────────────────────────────────
\begin{frame}{Scoring --- Normalization}
  \note{
    Before we can combine metrics into a final score, all raw values need to
    be on the same zero-to-one-hundred scale.  We use two approaches.

    For error-based metrics like MACE and onset MAE --- where lower is better
    --- we use piecewise-linear interpolation through a set of breakpoints.
    You pick a (x, score) pair for each perceptual threshold.  For example
    for MACE: 0 cents maps to 100, 25 cents maps to 88, 50 cents maps to 75,
    and 100 cents --- one full semitone --- maps to 50.  Values outside the
    range are clamped.  This lets us tune where the ``good'', ``acceptable'',
    and ``poor'' thresholds are for each metric independently.

    The breakpoints table on the right shows the values hard-coded in the
    source for all six error metrics.  Notice that the phoneme boundary
    breakpoints are stricter than the note onset breakpoints, reflecting that
    phoneme timing is a finer-grained task.

    For accuracy fractions like Pitch Accuracy and Timing Accuracy --- where
    the value is already between 0 and 1 and higher is better --- we simply
    multiply by 100.  No curve needed.

    There is also a Gaussian penalty function available in the codebase, but
    it is not called in the default scoring configuration.
  }
  \begin{columns}[T]
    \column{0.48\textwidth}
      \begin{block}{Piecewise-Linear (error metrics)}
        \[
          S_\text{pw}(x)
            = s_i + \frac{s_{i+1}-s_i}{x_{i+1}-x_i}(x - x_i),
            \quad x_i \le x \le x_{i+1}
        \]
        \footnotesize Clamped to $[0,100]$. Used for MACE, RMSE, onset MAE, IOI MAE,
        RelDurErr, PhonBndErr.
      \end{block}
      \medskip
      \begin{block}{Direct Scaling (accuracy fractions)}
        \[
          S_\text{acc}(a) = a \times 100
        \]
        \footnotesize Used for PitchAcc, TimingAcc, WordAcc, OverlapAcc, LabelMatch.
      \end{block}

    \column{0.48\textwidth}
      \begin{exampleblock}{Breakpoints — from source code}
        \footnotesize
        \renewcommand{\arraystretch}{1.2}
        \begin{tabular}{@{}lcc@{}}
          \toprule
          Metric & $x$ values & $s$ values \\
          \midrule
          MACE (¢)         & 0, 25, 50, 100, 200   & 100, 88, 75, 50, 0 \\
          Pitch RMSE (¢)   & 0, 25, 50, 100, 200   & 100, 88, 72, 45, 0 \\
          Onset MAE (ms)   & 0, 25, 50, 100, 200   & 100, 88, 75, 50, 0 \\
          IOI MAE (ms)     & 0, 30, 60, 120, 240   & 100, 88, 75, 50, 0 \\
          RelDurErr        & 0, .1, .2, .5, 1.0    & 100, 90, 75, 50, 0 \\
          PhonBndErr (ms)  & 0, 15, 30, 60, 120    & 100, 88, 75, 50, 0 \\
          \bottomrule
        \end{tabular}
      \end{exampleblock}
    \end{columns}
\end{frame}

%% ─── 9. Scoring — Category Scores ────────────────────────────────────────────
\begin{frame}{Scoring --- Category Scores}
  \note{
    Each of the four category scores is computed using the same
    confidence-weighted aggregation formula shown at the top.  We take a
    weighted sum of component scores, but each weight is multiplied by a
    confidence factor c-i.  Confidence is the minimum of 1 and N divided by
    3, so you need at least 3 matched events to reach full confidence.
    Components with no data contribute zero and effectively drop out of the
    denominator.

    The four component weight tables show what goes into each category.
    Pitch is 50 percent accuracy, 30 percent MACE, and 20 percent RMSE.
    Timing is 50 percent timing accuracy, 30 percent onset MAE, and 20 percent
    rhythm stability.  Duration is 60 percent relative error, 20 percent ratio
    deviation, and 20 percent phrase consistency.  Lyric is 35 percent word
    accuracy, 25 percent overlap, 25 percent label match, and 15 percent
    phoneme boundary timing.

    The two formulas at the bottom show how the two compound components are
    calculated.  Rhythm Stability is the average of the IOI-MAE score and the
    onset standard-deviation score --- it captures both inter-note interval
    consistency and the tightness of timing errors.  Phrase Consistency
    averages the duration standard-deviation score and the relative duration
    error score.  Either term is used alone if the other is unavailable.
  }
  \begin{block}{Confidence-Weighted Aggregation (all four categories)}
    \[
      S_\text{cat}
        = \frac{\displaystyle\sum_i w_i\,c_i\,s_i}
               {\displaystyle\sum_i w_i\,c_i},
      \qquad
      c_i = \min\!\Bigl(1,\,\tfrac{N}{3}\Bigr)
    \]
    \footnotesize $w_i$ = nominal weight;\; $c_i$ = confidence from data volume;\;
    $s_i\in[0,100]$ = component score.
  \end{block}

  \medskip
  \begin{columns}[T]
    \column{0.22\textwidth}
      \begin{exampleblock}{\footnotesize Pitch}
        \footnotesize
        \begin{tabular}{@{}lr@{}}
          PitchAcc & 0.50 \\
          MACE     & 0.30 \\
          RMSE     & 0.20 \\
        \end{tabular}
      \end{exampleblock}
    \column{0.22\textwidth}
      \begin{exampleblock}{\footnotesize Timing}
        \footnotesize
        \begin{tabular}{@{}lr@{}}
          TimingAcc & 0.50 \\
          Onset MAE & 0.30 \\
          IOI+std   & 0.20 \\
        \end{tabular}
      \end{exampleblock}
    \column{0.22\textwidth}
      \begin{exampleblock}{\footnotesize Duration}
        \footnotesize
        \begin{tabular}{@{}lr@{}}
          RelDurErr & 0.60 \\
          RatioDev  & 0.20 \\
          PhraseCon & 0.20 \\
        \end{tabular}
      \end{exampleblock}
    \column{0.22\textwidth}
      \begin{exampleblock}{\footnotesize Lyric}
        \footnotesize
        \begin{tabular}{@{}lr@{}}
          WordAcc     & 0.35 \\
          OverlapAcc  & 0.25 \\
          LabelMatch  & 0.25 \\
          PhonBndErr  & 0.15 \\
        \end{tabular}
      \end{exampleblock}
  \end{columns}

  \medskip
  \begin{columns}[T]
    \column{0.48\textwidth}
      \begin{block}{\footnotesize Rhythm Stability (Timing, $w=0.20$)}
        \footnotesize
        \[
          S_\text{rhythm}
            = \tfrac{1}{2}\bigl(S_\text{pw}(\mathrm{IOI\text{-}MAE}) + S_\text{pw}(\sigma_\text{on})\bigr)
        \]
      \end{block}
    \column{0.48\textwidth}
      \begin{block}{\footnotesize Phrase Consistency (Duration, $w=0.20$)}
        \footnotesize
        \[
          S_\text{phrase}
            = \tfrac{1}{2}\bigl(S_\text{pw}(\sigma_\text{dur}) + S_\text{pw}(\mathrm{RelDurErr})\bigr)
        \]
      \end{block}
  \end{columns}
\end{frame}

%% ─── 10. Scoring — Final Score & Interpretation ──────────────────────────────
\begin{frame}{Scoring --- Final Score \& Interpretation}
  \note{
    The final coaching score is a weighted combination of the four category
    scores: 40 percent pitch, 30 percent timing, 15 percent duration, and 15
    percent lyric.  Pitch has the highest weight because intonation is the most
    fundamental aspect of singing quality.

    At runtime the formula is actually computed using the confidence-adjusted
    version shown below the boxed equation, where each category weight is also
    multiplied by its confidence factor.  This means a category with very few
    evaluated notes contributes proportionally less to the overall score.

    The interpretation thresholds come directly from the configuration file.
    90 and above is excellent.  75 to 89 is good.  55 to 74 is fair.  Below
    55 needs work.

    All feedback messages are looked up from a static rule table keyed by
    category and performance level.  There are no language models involved.
    The output is completely deterministic --- the same audio file will always
    produce the same score and the same feedback.
  }
  \begin{block}{Final Score (\texttt{scoring/performance\_scoring.py})}
    \[
      \boxed{S_\text{final}
        = 0.40\,S_\text{pitch}
        + 0.30\,S_\text{timing}
        + 0.15\,S_\text{dur}
        + 0.15\,S_\text{lyric}}
    \]
    \[
      S_\text{final}
        = \frac{\sum_k W_k\,C_k\,S_k}{\sum_k W_k\,C_k},\quad
      C_k = \min\!\bigl(1,\,N_k/3\bigr)
      \quad\text{(confidence-adjusted)}
    \]
  \end{block}

  \medskip
  \begin{columns}[T]
    \column{0.50\textwidth}
      \begin{block}{Score Levels (\texttt{scoring/interpretation.py})}
        \small
        \renewcommand{\arraystretch}{1.2}
        \begin{tabular}{@{}clp{4.2cm}@{}}
          \toprule
          Score    & Level & Meaning \\
          \midrule
          $\ge 90$ & excellent  & Performance-ready \\
          $\ge 75$ & good       & A few targeted corrections \\
          $\ge 55$ & fair       & Focused practice needed \\
          $< 55$   & needs\_work & Coaching before full-song practice \\
          \bottomrule
        \end{tabular}
      \end{block}

    \column{0.46\textwidth}
      \begin{alertblock}{Rule-Based Feedback (no LLMs)}
        \small
        Strengths and weaknesses are looked up from a \textbf{static table}
        keyed by (category, level) — output is fully deterministic.
        \medskip\par
        \footnotesize
        \textit{Example strength:} ``Strong pitch intonation and accuracy''\\
        \textit{Example weakness:} ``Significant timing inconsistencies detected''
      \end{alertblock}
  \end{columns}
\end{frame}

%% ─── 11. Model — Pitch Extraction ────────────────────────────────────────────
\begin{frame}{Model --- Pitch Extraction (\texttt{models/pitch/})}
  \note{
    The pitch extraction pipeline has five stages.

    First, the WAV is loaded at 16 kHz.

    Second, WebRTC VAD runs with aggressiveness 2, 20-millisecond frames, and
    a median smoothing window of 5, producing a voiced/unvoiced mask at the VAD
    frame rate.  If WebRTC is not installed, an energy-based fallback activates
    at minus 40 dBFS.

    Third, the actual pitch is estimated.  The primary backend is torchcrepe, a
    CREPE-style CNN with Viterbi decoding, a periodicity threshold of 0.21, and
    a silence threshold of minus 60 dBFS.  The fallback is librosa's pYIN
    algorithm with its own frequency range of 65 to 2093 Hz.

    Fourth, the VAD mask is resampled to the pitch frame rate using
    nearest-neighbour interpolation.

    Fifth, fuse\_vad\_and\_pitch combines everything in five sub-steps.  It
    intersects the pitch model's voiced flags with the VAD mask.  It zeros
    unvoiced frames.  It linearly interpolates short gaps up to 10 frames and
    marks those frames as voiced.  It applies a median filter of kernel size 5
    to voiced frames.  Finally it applies Gaussian smoothing with sigma 1.

    The output is timestamps, F0 in Hz, and a voiced boolean mask, all at 100
    frames per second.
  }
  \begin{columns}[T]
    \column{0.50\textwidth}
      \begin{block}{Pipeline (\texttt{pipeline.py})}
        \small
        \renewcommand{\arraystretch}{1.2}
        \begin{tabular}{@{}lp{5.6cm}@{}}
          \textbf{1.\ WAV}   & 16\,kHz, mono, float32 \\
          \textbf{2.\ VAD}   & WebRTC VAD: aggressiveness=2, frame=20\,ms,
                               smoothing window=5 \\
          \textbf{3.\ Pitch} & torchcrepe (CREPE CNN + Viterbi,
                               period.\ thresh.=0.21, silence=$-60$\,dBFS)\\
                             & or librosa.pyin (fmin=65, fmax=2093\,Hz) \\
          \textbf{4.\ Align} & Resample VAD mask to pitch frame rate \\
          \textbf{5.\ Fuse}  & See right column \\
        \end{tabular}
      \end{block}
      \vspace{4pt}
      \begin{exampleblock}{Output}
        \small $(\,\text{timestamps},\; \fh_t\text{ (Hz)},\; \vh_t\text{ (bool)}\,)$
        at 100\,fps \; (hop = 160 samples at 16\,kHz)
      \end{exampleblock}

    \column{0.46\textwidth}
      \begin{block}{\texttt{fuse\_vad\_and\_pitch()} (\texttt{fusion.py})}
        \small
        \begin{enumerate}\setlength\itemsep{3pt}
          \item $\vh_\text{final} \leftarrow \vh_\text{pitch} \wedge \vh_\text{vad}$
          \item Zero unvoiced:\; $\fh_t \leftarrow 0$ if $\neg\vh_\text{final},t$
          \item Interpolate gaps (linear, max 10 frames); update $\vh_\text{final}$
          \item Median filter voiced frames (kernel = 5)
          \item Gaussian smooth voiced frames ($\sigma = 1.0$)
        \end{enumerate}
      \end{block}
  \end{columns}
\end{frame}

%% ─── 12. Model — Onset/Offset Detection ─────────────────────────────────────
\begin{frame}{Model --- Onset/Offset Detection (\texttt{model\_v5.py})}
  \note{
    The onset and offset detection model is based on wav2vec2-base from
    Facebook.  The raw waveform goes into a frozen 7-layer CNN feature
    extractor with a total stride of 320 samples, giving a 50 Hz frame rate.
    The features then pass through a 12-layer transformer encoder with hidden
    size 768, which is fine-tuned during training.  After dropout at p equals
    0.1, a single linear layer maps the 768-dimensional hidden state to 3
    logits per frame.

    Those three logits correspond to onset probability, offset probability,
    and active or voiced probability.

    At inference time, each logit is passed through a sigmoid to get a
    probability curve.  We then run scipy's find\_peaks on the onset and offset
    curves with a threshold of 0.3 and a minimum peak distance of 3 frames,
    which is about 60 milliseconds.  Onsets and offsets are then paired
    greedily: each onset is matched with the next available offset that comes
    after it.

    The model outputs at 50 Hz but is resampled to the canonical 100 Hz grid
    in the fusion step.

    The registry in the codebase automatically selects this Wav2Vec2 model
    when the checkpoint file ends in .ckpt, which is what pipeline.yaml
    specifies.
  }
  \textbf{Wav2Vec2OnsetDetector} — \texttt{facebook/wav2vec2-base} with a
  3-class linear head.

  \begin{columns}[T]
    \column{0.50\textwidth}
      \begin{block}{Architecture}
        \small
        \renewcommand{\arraystretch}{1.2}
        \begin{tabular}{@{}lp{5.6cm}@{}}
          Input       & Waveform $\mathbf{x}\in\mathbb{R}^{B\times T}$ at 16\,kHz \\
          CNN Extr.   & 7-layer Conv1d; total stride 320;
                        \textbf{frozen} \\
          Transformer & 12-layer encoder, $d=768$; fine-tuned \\
          Dropout     & $p=0.1$ \\
          Head        & Linear$(768\!\to\!3)$: logits $\ell^{(0)}_t,\ell^{(1)}_t,\ell^{(2)}_t$ \\
        \end{tabular}
        \vspace{2pt}
        \footnotesize
        $\ell^{(0)}$ = onset, $\ell^{(1)}$ = offset, $\ell^{(2)}$ = active (voiced)
      \end{block}

    \column{0.46\textwidth}
      \begin{block}{Inference}
        \[
          p_t^{(k)} = \sigma\!\left(\ell_t^{(k)}\right) \in [0,1]
        \]
        \[
          \sh_n = \mathrm{PeakPick}\!\left(p_t^{(0)},\;\theta_\text{on}{=}0.3\right)
        \]
        \[
          \eh_n = \mathrm{PeakPick}\!\left(p_t^{(1)},\;\theta_\text{off}{=}0.3\right)
        \]
        \footnotesize Min distance = 3 frames\,$\approx$\,60\,ms;
        onset--offset paired greedily.
      \end{block}
      \vspace{4pt}
      \begin{exampleblock}{Frame Rate}
        \small $16000 / 320 = 50$\,Hz (20\,ms/frame)\\
        Resampled to 100\,fps canonical grid.
      \end{exampleblock}
  \end{columns}
\end{frame}

%% ─── 13. Model — Phoneme Detection ──────────────────────────────────────────
\begin{frame}{Model --- Phoneme Boundary Detection (\texttt{phoneme\_model.py})}
  \note{
    The phoneme boundary model uses the pretrained wav2vec2-lv-60-espeak-cv-ft
    model from Facebook.  This was trained on 60 languages using an eSpeak IPA
    vocabulary.  Importantly, unlike the onset/offset model, this one is loaded
    purely in eval mode and used as-is for inference --- we are not fine-tuning
    any layers.

    The pipeline works as follows.  The audio goes through the Wav2Vec2 feature
    extractor CNN, then through the transformer encoder to produce per-frame
    logits over the phoneme vocabulary.  CTC alignment then takes argmax per
    frame, collapses consecutive repeated tokens, and removes blank tokens,
    leaving a sequence of timed phoneme segments.

    Two post-processing steps follow.  First, any segment longer than 300
    milliseconds is split at internal transition points where an alternative
    phoneme has a probability above 0.25.  This handles long sustained vowels.
    Second, a blank-region scan goes through the gaps between segments and
    inserts any phoneme whose peak probability exceeds 0.05 --- this recovers
    phonemes that CTC blank dominance tends to suppress in singing.

    For scoring, we match these predicted segments against the TextGrid
    reference boundaries and compute the three lyric metrics shown on the right.
  }
  \textbf{Pretrained:} \texttt{facebook/wav2vec2-lv-60-espeak-cv-ft}
  — loaded in eval mode (inference only, no frozen-layer training).

  \begin{columns}[T]
    \column{0.50\textwidth}
      \begin{block}{CTC Pipeline}
        \small
        \renewcommand{\arraystretch}{1.2}
        \begin{tabular}{@{}lp{5.5cm}@{}}
          \textbf{1.} Input  & WAV at 16\,kHz \\
          \textbf{2.} Feats  & Wav2Vec2FeatureExtractor (CNN, eval mode) \\
          \textbf{3.} Encode & Wav2Vec2 Transformer;
                               logits $\in\mathbb{R}^{T\times V}$ \\
          \textbf{4.} CTC    & Argmax; collapse repeated; remove blanks \\
          \textbf{5.} Post   & (i) split segments $>300$\,ms\\
                             & (ii) blank-region phoneme recovery \\
          \textbf{6.} Out    & List$[\,(\ah_p,\bh_p,\text{label},\text{conf})\,]$ \\
        \end{tabular}
      \end{block}

    \column{0.46\textwidth}
      \begin{block}{Scoring Connection}
        \[
          \mathrm{PhonBndErr}
            = \frac{1}{P}\sum_{p=1}^{P}|\ah_p - a_p|\times 1000
        \]
        \footnotesize Onset deviation only (not average of start and end).
        \medskip\par
        \small Additional outputs per matched pair:
        \begin{itemize}\setlength\itemsep{1pt}\footnotesize
          \item overlap fraction $\ge 0.5$ (OverlapAcc)
          \item label match flag $\hat{l}_p = l_p$ (LabelMatch)
        \end{itemize}
      \end{block}
  \end{columns}
\end{frame}

%% ─── 14. Results — Pitch Model ───────────────────────────────────────────────

%% ─── 15. Results — Onset/Offset ─────────────────────────────────────────────
\begin{frame}{Results --- Onset/Offset Detection (Wav2Vec2)}
  \note{
    This slide shows how the Wav2Vec2 onset/offset model's three outputs
    connect to the downstream scoring metrics.

    The onset logit produces S-hat-n, the predicted note start times.  These
    feed directly into onset error, timing accuracy, and IOI-MAE.

    The offset logit produces E-hat-n, the predicted note end times.  These
    feed into offset error, duration error, relative duration error, and
    duration ratio.

    The active logit provides a per-frame voiced activity mask that supplements
    the VAD signal during the note event construction phase.

    The peak-picking parameters are straightforward: a sigmoid threshold of 0.3
    for both onset and offset, a minimum inter-peak distance of 3 frames which
    is about 60 milliseconds, and greedy pairing where each onset is matched
    to the next available offset.

    The model runs at 50 Hz because of the Wav2Vec2 CNN stride of 320 samples
    at 16 kHz.  It is resampled to the canonical 100 Hz grid in the fusion
    step before any metric computation.
  }
  \begin{columns}[T]
    \column{0.55\textwidth}
      \begin{block}{Model Outputs and Downstream Metrics}
        \small
        \renewcommand{\arraystretch}{1.2}
        \begin{tabular}{@{}lp{5.6cm}@{}}
          \toprule
          Output & Used to compute \\
          \midrule
          $\sh_n$ from $\ell^{(0)}$ & OnsetError, TimingAcc, IOI-MAE \\
          $\eh_n$ from $\ell^{(1)}$ & OffsetError, DurErr, RelDurErr, DurRatio \\
          active from $\ell^{(2)}$ & voiced-frame gating in event construction \\
          \bottomrule
        \end{tabular}
      \end{block}
      \vspace{4pt}
      \begin{alertblock}{Frame Rate}
        \small
        CNN stride $= 320$ samples @ 16\,kHz $\Rightarrow$ 50\,Hz (20\,ms/frame)\\
        Resampled to 100\,fps canonical grid.
      \end{alertblock}

    \column{0.41\textwidth}
      \begin{block}{Peak-Picking Parameters}
        \small
        \renewcommand{\arraystretch}{1.2}
        \begin{tabular}{@{}ll@{}}
          \toprule
          Parameter & Value \\
          \midrule
          $\theta_\text{onset}$  & 0.3 \\
          $\theta_\text{offset}$ & 0.3 \\
          Min distance  & 3 frames ($\approx60$\,ms) \\
          Pairing       & Greedy, next available offset \\
          \bottomrule
        \end{tabular}
      \end{block}
  \end{columns}
\end{frame}

%% ─── 16. Results — Note-Level Pitch ─────────────────────────────────────────
\begin{frame}{Results --- Note-Level Pitch Accuracy (Pilot)}
  \note{
    This is a pilot evaluation on a 7-note excerpt from the song North Wind
    Meets the Sea, using the torchcrepe backend with a 50-cent tolerance.

    The overall note-level pitch accuracy is 6 out of 7 notes, or 85.7 percent.

    Looking at the table, note 2 on the lyric ``the'' is the only failure at
    minus 152.8 cents --- about one and a half semitones flat.  This is most
    likely because it is a very short note at only 178 milliseconds with just
    18 voiced frames, which is not enough data for a reliable median pitch
    estimate.

    All 6 correct notes are within 45 cents.  Note 6 at 44.4 cents is the
    closest to the 50-cent boundary.

    There is also a mild systematic bias in the results: the model tends to be
    flat on the longer sustained notes --- notes 1 and 4 --- and slightly sharp
    on the shorter faster notes.  This is worth investigating further as we
    gather more data.
  }
  \textbf{Setup:} 7-note excerpt $\cdot$ torchcrepe $\cdot$ $\tau=50$\
  \[
    \mathrm{NotePitchAcc}_{50} = 6/7 \approx 85.7\%
  \]
  \begin{center}
    \small
    \begin{tabular}{@{}clccccc@{}}
      \toprule
      Note & Lyric & Ref (MIDI) & Ref (Hz) & Det.\ (MIDI)
           & $\bar{c}_n$ (¢) & Correct \\
      \midrule
      1 & \textit{where} & C4 (60.0)  & 261.6 & 59.74 & $-26.1$  & $\checkmark$ \\
      2 & \textit{the}   & D4 (62.0)  & 293.7 & 60.47 & $-152.8$ & $\times$ \\
      3 & \textit{north} & E4 (64.0)  & 329.6 & 64.08 & $+8.2$   & $\checkmark$ \\
      4 & \textit{wind}  & Bb3 (58.0) & 233.1 & 57.63 & $-36.8$  & $\checkmark$ \\
      5 & \textit{meets} & C4 (60.0)  & 261.6 & 60.32 & $+32.5$  & $\checkmark$ \\
      6 & \textit{the}   & A3 (57.0)  & 220.0 & 57.44 & $+44.4$  & $\checkmark$ \\
      7 & \textit{sea}   & D4 (62.0)  & 293.7 & 62.24 & $+23.7$  & $\checkmark$ \\
      \bottomrule
    \end{tabular}
  \end{center}
  \begin{itemize}\setlength\itemsep{1pt}\footnotesize
    \item Note 2: $-152.8$\,¢ on a 178\,ms note with only 18 voiced frames
    \item All 6 correct notes within $\pm45$\,¢; Note 6 at 44.4\,¢ closest to the $\pm50$\,¢ boundary
    \item Mild bias: flat on sustained notes (1, 4); sharp on shorter notes (3, 5, 6, 7)
  \end{itemize}
\end{frame}

%% ─── 17. Integrated Pipeline ─────────────────────────────────────────────────
\begin{frame}{Integrated System --- VocalCoach Pipeline}
  \note{
    This final slide shows the complete integrated pipeline.

    Phase 1 loads the audio at 16 kHz with normalisation.

    Phases 2, 3, and 4 run the three models.  Phase 2 runs the pitch pipeline
    via torchcrepe and produces F-hat-t and V-hat-t at 100 fps.  Phase 3 runs
    the phoneme model via Wav2Vec2 CTC and produces a list of phoneme segments
    with timestamps.  Phase 4 runs the Wav2Vec2 onset/offset detector and
    produces S-hat-n and E-hat-n.

    Phase 5 is the feature fusion step.  All model outputs are merged onto the
    canonical 100 Hz frame grid, and then structured note events and lyric
    events are built into the FusedPerformanceRepresentation.

    Phase 6 parses the MusicXML and TextGrid reference files, then runs the
    greedy overlap-based alignment to match predicted events against reference
    events, producing the AlignmentResult.

    Phase 7 computes all the metrics we discussed, passes them through the
    scoring engine, and generates the interpretation summary.

    The entry point is a single call to UnifiedInferencePipeline.predict with
    the audio path, MusicXML path, TextGrid path, and compute\_scores equals
    True.

    To close: the key design principle is that the neural models only estimate
    measurable physical parameters.  All the coaching intelligence lives in
    deterministic equations, making the system fully interpretable and
    auditable.  Thank you.
  }
  \small
  \renewcommand{\arraystretch}{1.2}
  \begin{tabular}{@{}clp{10cm}@{}}
    \toprule
    Ph. & Module & Output \\
    \midrule
    1 & \texttt{utils/audio.py}
      & Load audio; 16\,kHz mono; canonical 10\,ms grid (hop=160) \\
    2 & \texttt{models/pitch/pipeline.py}
      & torchcrepe $\to$ $\fh_t$ (Hz), $\vh_t$ (bool) @ 100\,fps \\
    3 & \texttt{models/phoneme/phoneme\_model.py}
      & Wav2Vec2+CTC $\to$ List$[\,(\ah_p, \bh_p)\,]$ \\
    4 & \texttt{models/onset\_offset/model\_v5.py}
      & Wav2Vec2 $\to$ logits $\to$ $\sh_n, \eh_n$ \\
    5 & \texttt{fusion/}
      & Merge to 100\,fps; build NoteEvents, LyricEvents
        $\to$ FusedPerformanceRepresentation \\
    6 & \texttt{reference/} + \texttt{alignment/}
      & MusicXML+TextGrid $\to$ ReferenceRepresentation;
        greedy overlap $\to$ AlignmentResult \\
    7 & \texttt{metrics/} + \texttt{scoring/}
      & Metrics report $\to$ Score report $\to$ InterpretationSummary \\
    \bottomrule
  \end{tabular}

  \medskip
  \begin{alertblock}{Key Design Principle}
    \small
    Models estimate $\fh_t, \vh_t, \sh_n, \eh_n, \ah_p, \bh_p$.
    All coaching scores are computed by \textbf{deterministic equations}
    applied to those parameters — models never directly predict a score.
  \end{alertblock}
\end{frame}

\end{document}
